\section{Kết luận & Hướng Phát triển}
\label{sec:conclusion}

\subsection{Tổng kết Kết quả Đạt được}

Báo cáo bài thực hành 07 đã trình bày toàn diện giải pháp thiết kế, cài đặt và đánh giá thực nghiệm hệ thống phân vùng bản đồ ngữ nghĩa và điều hướng tự trị bằng giọng nói cho robot TurtleBot 4 trong môi trường Ignition Gazebo và ROS 2 Humble. Các kết quả nổi bật bao gồm:

\begin{enumerate}[label=\textbf{\arabic*.}]
    \item \textbf{Phân vùng Không gian Ngữ nghĩa Tự động:} Làm chủ và hiện thực hóa thành công giải pháp phân vùng bản đồ lưới chiếm chỗ $14\,\text{m} \times 11\,\text{m}$ bằng giải thuật Marker-controlled Watershed và biến đổi khoảng cách Euclidean. Hệ thống tự động nhận diện và trích xuất hoàn hảo 5 phòng chức năng với chỉ số IoU trung bình đạt $\mathbf{98.0\%}$, tự động sinh các điểm đón an toàn (\textit{safe waypoints}) tại tâm cửa mở rộng $1.3\,\text{m}$.
    
    \item \textbf{Giải pháp Xử lý Âm thanh Chống Méo Quá Ngưỡng:} Đề xuất và cài đặt thành công chuỗi tiền xử lý tín hiệu âm thanh 5 tầng gồm lọc thông cao 80\,Hz, khử trôi DC, ngắt khoảng lặng thông minh (\textit{Smart Noise Gate / VAD}), nắn méo bão hòa mềm $\tanh$ và tự động cân bằng âm lượng thích nghi (Auto AGC đưa đỉnh âm về $-9\,\text{dBFS}$). Giải pháp này đã \textbf{khắc phục triệt để lỗi nhận diện sai khi tăng âm lượng micro lên $100\%$ để quay video màn hình}, nâng độ chính xác nhận dạng câu lệnh từ $14.0\%$ lên $\mathbf{98.8\%}$.
    
    \item \textbf{Kiến trúc Nhận dạng Giọng nói Đa tầng Tin cậy:} Đảo chiều thứ bậc ưu tiên, đưa Google Cloud Deep Learning Speech AI làm bộ nhận dạng chính với độ chuẩn xác tiếng Việt vượt trội, kết hợp mô hình offline Vosk có màng lọc từ khóa điều hướng (\textit{Domain Keyword Guard}), loại bỏ hoàn toàn hiện tượng nhận nhầm tạp âm quạt tản nhiệt thành các từ ngữ ảo giác.
    
    \item \textbf{Bộ Hiểu Ngôn ngữ Tự nhiên & Lập kế hoạch Đa chặng:} Xây dựng bộ NLU Parser vượt qua trọn vẹn $\mathbf{104 / 104}$ ca kiểm thử tự động ($100\%$), hỗ trợ linh hoạt các mệnh lệnh đơn chặng, dừng khẩn cấp và các chuỗi hành trình phức tạp đa chặng liên phòng (*"tới khu E rồi tới văn phòng"*).
    
    \item \textbf{Điều hướng Mượt mà & Giao diện Trực quan Thời gian thực:} Tối ưu hóa bộ điều khiển DWB và bán kính lạm phát chi phí, giúp robot luồn lách êm ái qua cửa hẹp với tỷ lệ thành công $\mathbf{100\%}$ không đơ dừng; tích hợp cùng giao diện Studio PyQt5 hiện đại, cung cấp trải nghiệm điều khiển HRI trực quan và chuyên nghiệp.
\end{enumerate}

\subsection{Bài học Kinh nghiệm & Hướng Phát triển}

Quá trình triển khai đề tài mang lại những bài học kinh nghiệm sâu sắc:
\begin{itemize}
    \item \textbf{Sự phụ thuộc giữa phần cứng thu âm và mô hình AI:} Trong các hệ thống HRI thực tế, chất lượng tín hiệu âm học đầu vào đóng vai trò quyết định. Các giải pháp lọc nhiễu số và AGC là mắt xích bắt buộc để cầu nối giữa giọng nói con người và mô hình nhận diện không bị đứt gãy do biến thiên âm lượng phần cứng.
    \item \textbf{Gắn kết ngữ nghĩa với điều khiển vật lý:} Phân vùng bản đồ không chỉ dừng lại ở xử lý ảnh mà phải tích hợp chặt chẽ với động học robot. Vị trí waypoint và vùng lạm phát costmap phải được tính toán dung sai cẩn trọng để đảm bảo khả năng lập quỹ đạo khả thi trong Nav2.
\end{itemize}

Trong tương lai, hệ thống có thể mở rộng theo các hướng:
\begin{enumerate}
    \item Tích hợp mô hình ngôn ngữ lớn chạy trên biên (\textit{Edge LLM} như Llama 3 / Gemma) để xử lý các câu lệnh mang tính suy luận ngữ cảnh cao (*"tìm cho tôi một nơi yên tĩnh để học bài"* $\rightarrow$ tự động điều hướng đến Khu D: Thư viện).
    \item Kết hợp thị giác máy tính nhận diện khuôn mặt và cử chỉ tay để tạo thành hệ thống HRI đa phương thức hoàn chỉnh (\textit{Multimodal Voice-Vision Control}).
\end{enumerate}

\subsection{Thông tin Liên kết Báo cáo}

Toàn bộ mã nguồn, cấu hình mô phỏng, kịch bản kiểm thử và video ghi hình thực nghiệm được lưu trữ công khai tại các liên kết sau:
\begin{itemize}
    \item \textbf{Mã nguồn Dự án (GitHub):} \\
    \url{https://github.com/bangrbt/turtlebot4_lab/tree/assignments_4/src/turtlebot4_voice_navigation}
    \item \textbf{Video Thực nghiệm & Demo Hoạt động (Google Drive):} \\
    \url{https://drive.google.com/file/d/13ebzpQBoCgz06Z1LBtly1qYH1US6P3N8/view?usp=sharing}
\end{itemize}
